\documentclass[runningheads]{llncs}

\usepackage[T1]{fontenc}
\usepackage[utf8]{inputenc}
\usepackage{lmodern,microtype}
\usepackage{amsmath,amssymb,mathtools}
\usepackage{array,booktabs,enumitem}
\usepackage{bussproofs}
\usepackage{hyperref}
\hypersetup{colorlinks=true,linkcolor=blue,citecolor=blue,urlcolor=blue}
\usepackage{xcolor}

\input{../cyclic-cic/macros}

\newcommand{\ciu}{\approx_{\mathrm{ciu}}}
\newcommand{\tconf}[3]{#1 \vdash #2 : #3}

\title{Backlink Admissibility in Cyclic Proofs\\
       for the Calculus of Inductive Constructions}

\author{Gavin Mendel-Gleason}
\institute{Scidonia Limited\\
\email{gavin@scidonia.ai}}

\begin{document}
\maketitle

\begin{abstract}
Cyclic proofs close goals by backlinking to previously-seen configurations,
discovering induction principles post-hoc rather than prescribing them upfront.
We prove that backlinks are \emph{admissible} in the Calculus of Inductive
Constructions (CIC): every valid cyclic proof, presented as a finite directed
graph satisfying a budget-based progress condition, can be residualised to a
standard CIC term using only the native fixpoint binder ($\Fix{A}{t}$) for
recursion, and the residual is observationally indistinguishable from the root
goal.  Crucially, we work throughout with \emph{finitely presented} proof
structures; this finiteness is what allows the soundness argument to proceed by
well-founded induction entirely within CIC's inductive fragment, with no appeal
to coinduction or infinite proof trees.  The proof is mechanised in Rocq
and relies only on \emph{functional extensionality}---a standard, widely-accepted
axiom for CIC---with no other axioms or uses of \texttt{Admitted}.
\end{abstract}

\section{Introduction}

Cyclic proof theory~\cite{brotherston-2006-cyclic,brotherston-simpson-2011-sequent}
extends standard proof systems with backlink rules that allow a premise to
reference a previous goal, forming a directed cycle.  The cycle is sound if
it satisfies a global \emph{trace condition}: every cycle must contain a
\emph{progress event}---a case-split on a structurally smaller term.
Brotherston and Simpson prove that for first-order logic with inductive
definitions, backlinks are \emph{admissible}: every cyclic proof corresponds
to a standard proof with an explicit induction rule.

This paper extends their result to the Calculus of Inductive Constructions
(CIC)~\cite{coquand-huet-1988-coc,paulin-mohring-1993-inductive}, the type
theory underlying Rocq.  We prove that every rooted, ranked cyclic proof
can be residualised to a standard CIC term using only $\Fix{A}{t}$ for
recursion, and that the residual is observationally equivalent to the
root goal.

\paragraph{Finite presentations and the inductive fragment.}

A fundamental point separates our account from Brotherston's original
treatment~\cite{brotherston-2006-cyclic}.  Brotherston works with the
\emph{infinite regular tree} obtained by unrolling a cyclic proof, and proves
soundness by reasoning about infinite branches of that tree---an argument that
inherently requires some form of limit reasoning (K\"{o}nig's lemma, or a
semantic argument about infinite derivations).

We work throughout with the \emph{finitely presented} proof structure itself.
A cyclic proof is a \emph{finite directed graph} whose vertex set $V$ is
finite, decidable, and countable (concretely, $\mathbb{N}$), with each vertex
labelled by a typing judgement and its outgoing edges labelled by the rule
premises applied at that vertex.  This finiteness is not a restriction in
practice: every cyclic proof produced by a proof-search or supercompilation
procedure is finite by construction.

The finiteness of $V$ is what makes the soundness argument purely inductive.
The progress condition guarantees a strictly decreasing lexicographic measure
across every back-edge traversal in the \emph{finite} graph---so the
well-founded recursion terminates.  No unrolling to an infinite tree is
needed; no coinductive reasoning is required.  The proof term produced by
the \texttt{backlink\_admissible} theorem in Rocq is a fully normalising,
purely inductive term that Rocq's kernel verifies without any \texttt{cofix}
or classical axioms.

In short: \emph{the finite presentation of a cyclic proof induces an inductive
proof term}.  This is the mechanism by which cyclic proofs gain their expressive
power while remaining within CIC.

\paragraph{Future work: coinductive generalisation.}
An interesting direction is to extend this approach to proof systems where
the cyclic structure is not given in advance as a finite graph, but is instead
generated lazily or represented as an infinite but regular tree---as arises
naturally in proof-search over infinitely-branching or higher-order systems.
In such a setting the appropriate framework is \emph{coinductive}: the progress
condition would need to be reformulated as a coinductive invariant on the
infinite unrolling, and soundness would become a productivity argument in the
spirit of Coquand's coinductive proof theory or the global trace conditions
of Brotherston--Simpson.  We leave this generalisation as future work.

\section{Term calculus and typing}

\paragraph{Term syntax.}
Terms are given by the grammar:
\[
\begin{array}{rcl}
 t, A, B &::=& \Var{x} \mid \Sort{i} \mid \PiTy{A}{B} \mid \Lam{A}{t}
                \mid \App{t}{u} \mid \Fix{A}{t} \\
         &\mid& \Ind{I}{\vec{a}} \mid \Roll{I}{c}{\vec{a}}
                \mid \Case{I}{t}{C}{\vec{b}}
\end{array}
\]
where $x$ ranges over de~Bruijn indices, $i$ over universe levels,
$I$ identifies an inductive family, $c$ a constructor index,
$C$ a dependent motive, and $\vec{a},\vec{b}$ are lists of terms.

\paragraph{Typing (compact).}
A typing environment $\env$ stores inductive signatures.
A context $\ctx$ is a list of types.
The judgement $\hastype{\env}{\ctx}{t}{A}$ is standard for CIC with
inductives~\cite{coquand-huet-1988-coc,paulin-mohring-1993-inductive}.
The core rules are summarised in Figure~\ref{fig:typing}.

\begin{figure}[t]
\centering
\begin{minipage}{0.48\linewidth}
\begin{prooftree}
  \AxiomC{$\ctx(x)=A$}
  \RightLabel{\scriptsize (TyVar)}
  \UnaryInfC{$\hastype{\env}{\ctx}{\Var{x}}{A}$}
\end{prooftree}
\end{minipage}%
\begin{minipage}{0.48\linewidth}
\begin{prooftree}
  \AxiomC{}
  \RightLabel{\scriptsize (TySort)}
  \UnaryInfC{$\hastype{\env}{\ctx}{\Sort{i}}{\Sort{i+1}}$}
\end{prooftree}
\end{minipage}

\bigskip

\begin{minipage}{0.48\linewidth}
\begin{prooftree}
  \AxiomC{$\hastype{\env}{\ctx}{A}{\Sort{i}}$}
  \AxiomC{$\hastype{\env}{A{::}\ctx}{B}{\Sort{j}}$}
  \RightLabel{\scriptsize (TyPi)}
  \BinaryInfC{$\hastype{\env}{\ctx}{\PiTy{A}{B}}{\Sort{\max(i,j)}}$}
\end{prooftree}
\end{minipage}%
\begin{minipage}{0.48\linewidth}
\begin{prooftree}
  \AxiomC{$\hastype{\env}{\ctx}{A}{\Sort{i}}$}
  \AxiomC{$\hastype{\env}{A{::}\ctx}{t}{B}$}
  \RightLabel{\scriptsize (TyLam)}
  \BinaryInfC{$\hastype{\env}{\ctx}{\Lam{A}{t}}{\PiTy{A}{B}}$}
\end{prooftree}
\end{minipage}

\bigskip

\begin{minipage}{0.48\linewidth}
\begin{prooftree}
  \AxiomC{$\hastype{\env}{\ctx}{t}{\PiTy{A}{B}}$}
  \AxiomC{$\hastype{\env}{\ctx}{u}{A}$}
  \RightLabel{\scriptsize (TyApp)}
  \BinaryInfC{$\hastype{\env}{\ctx}{\App{t}{u}}{B[u/0]}$}
\end{prooftree}
\end{minipage}%
\begin{minipage}{0.48\linewidth}
\begin{prooftree}
  \AxiomC{$\hastype{\env}{\ctx}{A}{\Sort{i}}$}
  \AxiomC{$\hastype{\env}{A{::}\ctx}{t}{A\uparrow}$}
  \RightLabel{\scriptsize (TyFix)}
  \BinaryInfC{$\hastype{\env}{\ctx}{\Fix{A}{t}}{A}$}
\end{prooftree}
\end{minipage}

\bigskip

\begin{minipage}{0.48\linewidth}
\begin{prooftree}
  \AxiomC{$\env(I)=\Sigma_I$}
  \RightLabel{\scriptsize (TyInd)}
  \UnaryInfC{$\hastype{\env}{\ctx}{\Ind{I}{\vec{a}}}{\Sort{\ell}}$}
\end{prooftree}
\end{minipage}%
\begin{minipage}{0.48\linewidth}
\begin{prooftree}
  \AxiomC{$\hastype{\env}{\ctx}{a_i}{\tau_{c,i}}$ for each}
  \AxiomC{$\hastype{\env}{\ctx}{r_j}{\Ind{I}{}}$ for each re}
  \RightLabel{\scriptsize (TyRoll)}
  \BinaryInfC{$\hastype{\env}{\ctx}{\Roll{I}{c}{\vec{a}}}{\Ind{I}{\vec{\alpha}}}$}
\end{prooftree}
\end{minipage}

\bigskip

\begin{prooftree}
  \AxiomC{$\hastype{\env}{\ctx}{t}{\Ind{I}{\vec{\delta}}}$}
  \AxiomC{$\hastype{\env}{\ctx,\,x{:}\Ind{I}{\vec{\delta}}}{C}{\Sort{i}}$}
  \AxiomC{$\hastype{\env}{\ctx}{b_c}{\Pi\vec{\tau}_c.\,C[\Roll{I}{c}{\vec{y}}/x]}$  for each $c$}
  \RightLabel{\scriptsize (TyCase)}
  \TrinaryInfC{$\hastype{\env}{\ctx}{\Case{I}{t}{C}{\vec{b}}}{C[t/x]}$}
\end{prooftree}
\caption{Core typing rules of CIC with inductives.}
\label{fig:typing}
\end{figure}

\paragraph{Operational semantics.}
We use a call-by-name (CBN) weak-head reduction $\step$ defined by:
\[
\App{(\Lam{A}{t})}{u} \step t[u/0] \qquad
\Fix{A}{t} \step t[\Fix{A}{t}/0] \qquad
\Case{I}{(\Roll{I}{c}{\vec{a}})}{C}{\vec{b}} \step
  \mathit{apps}\;(\vec{b}_c)\,\vec{a}
\]
plus congruence under $\App{}{}$ and $\Case{}{}{}{}$ scrutinee position.
We write $\steps$ for the reflexive-transitive closure of $\step$, and
$t \Downarrow v$ for $t \steps v$ with $v$ a value (abstraction or
constructor). Multi-argument application $\mathit{apps}\;t\;[u_1,\dots,u_n]$
applies $t$ to each $u_i$ in turn.

\section{Cyclic sequent calculus}
\label{sec:calculus}

\paragraph{Finite configuration graphs.}
A \emph{configuration} is a typing judgement $\tconf{\Gamma}{t}{A}$.
A \emph{configuration graph} $b$ is a finite structure with vertex
set $V$ (implicitly $\{0,\dots,\mathit{next}-1\}$), a labelling
$\ell(b,\cdot):V\to\mathsf{Config}$, a successor map $\sigma(b,\cdot):V\to
\mathsf{list}(V)$, and a distinguished root $v_0\in V$.  In the
mechanisation the graph is a \texttt{cfg\_builder} record with finite
\texttt{gmap nat} tables; well-formedness (\texttt{cfg\_builder\_well\_formed})
ensures that every successor is itself a labelled vertex.

A graph is a \emph{rooted pre-proof} if every $v\in V$ satisfies one of the
local rules below (read bottom-up):

\paragraph{Asynchronous rules (deterministic).} Applied eagerly; each
reduces the goal to a single premise.

\medskip
\noindent
\begin{minipage}{0.48\linewidth}
\begin{prooftree}
  \AxiomC{$\tconf{\Gamma}{t[u/0]}{A}$}
  \RightLabel{\scriptsize ($\beta$)}
  \UnaryInfC{$\tconf{\Gamma}{\App{(\Lam{A}{t})}{u}}{A}$}
\end{prooftree}
\end{minipage}%
\begin{minipage}{0.48\linewidth}
\begin{prooftree}
  \AxiomC{$\tconf{\Gamma}{t[\Fix{A}{t}/0]}{A}$}
  \RightLabel{\scriptsize (fix)}
  \UnaryInfC{$\tconf{\Gamma}{\Fix{A}{t}}{A}$}
\end{prooftree}
\end{minipage}

\medskip
\noindent
\begin{minipage}{\linewidth}
\begin{prooftree}
  \AxiomC{$\tconf{\Gamma}{\mathit{apps}\;(\vec{b}_c)\;\vec{a}}{A}$}
  \RightLabel{\scriptsize ($\iota$)}
  \UnaryInfC{$\tconf{\Gamma}{\Case{I}{(\Roll{I}{c}{\vec{a}})}{C}{\vec{b}}}{A}$}
\end{prooftree}
\end{minipage}

\paragraph{Synchronous rule (choice point).} A case-split on a neutral
scrutinee~$x$ creates one premise per constructor $c$ of $I$, each
substituting the constructor pattern into the $c$-th branch:

\medskip
\noindent
\begin{minipage}{\linewidth}
\begin{prooftree}
  \AxiomC{$\tconf{\Gamma,\vec{y}_0}{b_0[\Roll{I}{0}{\vec{y}_0}/x]}{A}$}
  \AxiomC{$\cdots$}
  \AxiomC{$\tconf{\Gamma,\vec{y}_n}{b_n[\Roll{I}{n}{\vec{y}_n}/x]}{A}$}
  \RightLabel{\scriptsize (split)}
  \TrinaryInfC{$\tconf{\Gamma}{\Case{I}{x}{C}{\vec{b}}}{A}$}
\end{prooftree}
\end{minipage}
\medskip

\noindent A \textsc{split} vertex is a \emph{progress vertex}: it records a
case-split on a neutral scrutinee and is the locus of structural decrease
that underpins the trace condition.

\paragraph{Cyclic backlink (closure).}
A vertex with no premises may be closed not by an axiom, but by a
\emph{back-edge} to a previously-seen configuration under a
substitution~$\sigma$:

\medskip
\noindent
\begin{minipage}{\linewidth}
\begin{prooftree}
  \AxiomC{$\Delta \vdash \sigma :_s \Gamma$}
  \RightLabel{\scriptsize (back)}
  \UnaryInfC{$\tconf{\Delta}{t}{A}$}
\end{prooftree}
\end{minipage}
\medskip

\noindent where there exists a vertex $v' \in V$ such that
$\ell(b,v') = \tconf{\Gamma}{t_0}{A}$ and $t = t_0[\sigma]$ is the goal
under the substitution.

\paragraph{Global progress condition.}
A rooted pre-proof satisfies the progress condition if every directed
cycle contains at least one progress vertex.
In the mechanisation this is checked via a \emph{budget trace}:

\begin{definition}[Budget trace graph]
Given a configuration graph $G$ with progress predicate $\mathsf{prog}$
and initial budget $B = |V|$, the \emph{budget trace graph} has vertices
$V\times\{0,\dots,B\}$ and an edge $(v_1,k_1)\to(v_2,k_2)$ whenever
$v_1\to v_2$ is an edge in $G$ and:
\begin{itemize}
\item $v_1$ is a progress vertex and $k_2 = k_1 - 1$, or
\item $v_1$ is not a progress vertex and $k_2 = k_1$.
\end{itemize}
\end{definition}

\noindent The progress condition holds---written $\mathcal{C}(b,v_0)$---iff
the budget trace graph is acyclic.  In code:
$\texttt{trace\_condition\_ok}(b) \equiv
 \neg\,\texttt{has\_nonprogress\_cycle}(b)$.
Since the budget starts at $B=|V|$ and can decrease at most $|V|$ times,
acyclicity of the budget trace follows from the progress condition
(\texttt{trace\_rank\_monotone}, \texttt{trace\_rank\_strict\_on\_progress}).

\section{Residualisation}
\label{sec:residualise}

The backlink rule appears only in the proof graph.  A \emph{residualisation}
function $\mathcal{R}(b,v,d,\rho)$ reads back the graph into a CIC term by
traversing the successors of $v$, replacing backlinks with $\Fix{}{}$
binders and recording previously-seen vertices in a \emph{fix environment}
$\rho : \mathbb{N} \to \mathbb{N}$ (a finite map from graph vertices to
de~Bruijn depths).  Formally, $\mathcal{R}$ is defined by well-founded
recursion on the budget counter; the fuel parameter guarantees
termination in the mechanisation.

\begin{itemize}
\item \textbf{Leaf} ($v$ has no successors): returns $t$ where
  $\ell(b,v) = \tconf{\Gamma}{t}{A}$.
\item \textbf{Asynchronous} ($v$ has a single successor $w$):
  $\mathcal{R}(b,v,d,\rho) = D(\mathcal{R}(b,w,d,\rho))$ where $D$
  re-abstracts the driving step ($\beta$, fix-unfold, $\iota$).
\item \textbf{Split} ($v$ has multiple successors $\vec{w}$):
  each branch is residualised independently and combined into a
  $\Case{}{}{}{}$ application.
\item \textbf{Backlink} ($v$ is the target of a backlink, and
  $v \notin \mathrm{dom}(\rho)$): inserts a $\Fix{A}{t}$ binder where
  $t$ is the residual of the backlink's source vertex with the
  recursive self-reference variable at depth $d$.
\item \textbf{Already seen} ($v \in \mathrm{dom}(\rho)$): returns an
  application of the corresponding de~Bruijn variable.
\end{itemize}

The result is a well-typed CIC term using only standard constructors
($\lambda$, application, $\Fix{}{}$, $\Case{}{}{}{}$, $\Roll{}{}{}$).
No backlink constructs remain.

\paragraph{Why finiteness is essential.}
The residualisation terminates because the budget counter strictly
decreases on every back-edge and is bounded by $|V|$ which is finite.
If $V$ were infinite, the recursion would not terminate and the function
could not be given a type in CIC without coinduction.  The finite vertex
set is thus the direct reason the residualiser produces an \emph{inductive}
proof term.

\section{CIU equivalence}

The soundness guarantee is \emph{CIU equivalence}: for closed terms,
\[
  t \ciu u \;\equiv\;
  \bigl(\forall \sigma\, v.\; t[\sigma] \Downarrow v \Rightarrow u[\sigma] \Downarrow v\bigr)
  \;\wedge\;
  \bigl(\forall \sigma\, v.\; u[\sigma] \Downarrow v \Rightarrow t[\sigma] \Downarrow v\bigr)
\]
where $\sigma : \mathbb{N} \to \mathsf{tm}$ is an arbitrary de~Bruijn
substitution.  This is the open CIU equivalence lifted to all closing
substitutions.

Each local rule preserves CIU equivalence.  The asynchronous rules are
single reduction steps (CIU-preserving by \texttt{step\_ciu}); the split
rule applies CIU-congruence for neutral scrutinees; the back rule becomes a
$\Fix{}{}$ binder whose unfolding reproduces the source derivation.

\begin{theorem}[CIU soundness]
\label{thm:ciu}
If $\texttt{supercompile\_jTy\_tc}(\mathit{fuel},\Sigma,\Gamma,t,A) =
  \mathsf{Some}(v, b)$, then
\[
  t \ciu \mathcal{R}(b,v,0,\emptyset).
\]
\end{theorem}

\begin{proof}[Proof sketch]
\texttt{supercompile\_jTy\_tc} runs \texttt{supercompile\_jTy} to produce
a graph $b$ rooted at $v$, then checks \texttt{trace\_condition\_ok}$(b)$.
The CIU theorem (\texttt{supercompile\_ciu\_soundness\_untyped}) proceeds by
well-founded induction on the budget trace.  The well-foundedness uses
\texttt{lex\_lt\_wf}: the per-type lexicographic rank is well-founded and
every progress edge strictly decreases the budget
(\texttt{trace\_rank\_strict\_on\_progress}).  Critically, this induction is
over the \emph{finite} budget trace graph: there is no coinductive step,
no appeal to K\"{o}nig's lemma, and no reasoning about infinite unrollings.
The entire argument is \texttt{well\_founded\_induction lt\_wf} applied to
the budget counter, bounded by $|V|$.
\end{proof}

\section{Admissibility}

\begin{theorem}[Backlink admissibility]
\label{thm:admissible}
For any $\Gamma, t, A$ and fuel bounds $\mathit{fuel\_sc},
\mathit{fuel\_res}$, if
\[
  \texttt{supercompile\_jTy\_tc}(\mathit{fuel\_sc},\Sigma,\Gamma,t,A)
  = \mathsf{Some}(v, b),
\]
then there exists a standard CIC term $t_{\mathsf{std}}$ such that
\[
  t \ciu t_{\mathsf{std}}
  \quad\text{and}\quad
  t_{\mathsf{std}} = \mathcal{R}(b,v,0,\emptyset).
\]
\end{theorem}

\noindent In words: if the supercompiler produces a finite graph satisfying
the trace condition, the root term is CIU-equivalent to a standard CIC term
(using only $\lambda$, application, $\Fix{}{}$, $\Case{}{}{}{}$,
$\Roll{}{}{}$), with no cyclic backlink constructs.

\begin{proof}
Take $t_{\mathsf{std}} = \mathcal{R}(b,v,0,\emptyset)$.
By construction, $t_{\mathsf{std}}$ uses only standard CIC constructors.
CIU equivalence is Theorem~\ref{thm:ciu}.
\end{proof}

\noindent The mechanised proof (\texttt{BacklinkAdmissible.v},
theorem \texttt{backlink\_admissible}) is a four-line corollary: it
existentially introduces $t_{\mathsf{std}}$, splits the conjunction, and
applies \texttt{supercompile\_ciu\_soundness\_untyped} for CIU and
definitional unfolding for the equation.  The proof term is purely inductive
(no \texttt{cofix}, no \texttt{Admitted}).  The sole non-constructive
ingredient is \emph{functional extensionality}
($\forall f\,g,\,(\forall x,\,f(x)=g(x))\Rightarrow f=g$), which enters
via Autosubst's substitution lemmas and our typing substitution proofs.
This axiom is standard and widely accepted in CIC mechanisations; it does
not affect the computational content of the proof term.

\section{Discussion}

\paragraph{Finite presentation is essential, not incidental.}
The condition \texttt{supercompile\_jTy\_tc}$(\ldots) = \mathsf{Some}(v,b)$
encodes three things: (i)~the supercompiler terminates within the given fuel;
(ii)~it produces a closed finite graph; and (iii)~that graph satisfies the
trace condition.  All three are needed.  In particular, finiteness bounds
the budget $B = |V|$, makes the budget trace graph finite, and enables
well-founded recursion.

\paragraph{Relationship to Brotherston--Simpson.}
Brotherston and Simpson prove admissibility for first-order logic by
unrolling cyclic proofs to infinite regular trees and applying a semantic
argument about infinite derivations.  Our proof is structurally simpler: we
never leave the finite graph.  Instead of reasoning about infinite branches,
we reason about the budget counter, which decreases on every back-edge and
is bounded by $|V|$.  The finiteness of $V$ is what allows admissibility to
be expressed as \texttt{well\_founded\_induction} rather than as an ordinal
assignment or compactness argument.

\paragraph{The inductive/coinductive boundary.}
In CIC there is a sharp boundary between \emph{inductive} types
(well-founded recursion, terminating) and \emph{coinductive} types (guarded
corecursion, productive).  Cyclic proofs appear at first to belong on the
coinductive side: they contain cycles, and their natural models are infinite
regular trees.  The central insight of this paper is that when the cyclic
proof is \emph{finitely presented}---i.e.\ given as a concrete finite
graph---the progress condition provides enough structure to carry out the
soundness argument entirely on the inductive side.  The finite graph,
together with its progress condition, induces a well-founded measure, which
induces a terminating recursion, which induces a normalising proof term.

\paragraph{Acknowledgements.}
Thanks to James Brotherston for the question that motivated this paper.

\bibliographystyle{splncs04}
\bibliography{../cyclic-cic/refs}

\end{document}
